
\subsubsection{FPV Training Simulators}

Training Simulators are designed to imitate the operations and behaviors of a facility or process, so as to develop the skills or experiences of the operators who use them. The objective of the operator of the Training Simulator is to maximize his performance in achieving the task being simulated.\parencite{narayanasamy_distinguishing_2006}. Simulation-based learning as an active and immersive experience in a virtual learning environment that recreate a realistic and authentic real-life event or a set of learning situations to solve a problem.\parencite{dai_educational_2022}

In aviation and drone domains, such systems replicate flight dynamics, cockpit instrumentation, and environmental conditions of real aircraft. Flight training simulators have a long history, beginning around 1910 and evolving significantly during the World Wars. By the 1960s, they had become essential tools for pilot training, enabling the safe practice of complex scenarios. With advances in artificial intelligence and related technologies, modern simulators can now provide real-time feedback based on expert pilot performance. Today, they are widely used in both civilian and military contexts for training and mission planning \parencite{pappas_generic_2025}.

An FPV (First-Person View) training simulator is a specialized computer-based training system designed to teach and refine the technical skills required to operate FPV drones. Similar to traditional flight simulators, FPV simulators are widely used across entertainment, professional training, and research contexts. 

Existing literature on drone simulators has primarily focused on technical design or training effectiveness, often using research-oriented platforms such as Microsoft AirSim. However, widely used commercial FPV simulators (e.g., Liftoff and Uncrashed) remain under explored in pedagogical research. These commercial simulators typically emphasize high-fidelity environments and support advanced flight modes such as acrobatic mode for pilot practice. In contrast, research platforms such as AirSim and RotorS are widely used for machine learning and autonomous system development, but are not specifically tailored for practical pilot training \parencite{pappas_generic_2025}.

Despite their realism, FPV simulators also present several limitations. For example, while they aim to replicate real-world physics, users often report a steep learning curve and highly unforgiving flight dynamics, which may not always support beginner training effectively \parencite{guide_best_2026}. Additionally, hardware compatibility issues can arise, as not all radio transmitters are universally supported. The lack of standardization in flight controller architectures, along with the prevalence of proprietary and closed-source systems in UAV platforms, further complicates integration and usability \parencite{ebeid_survey_2018}.